%% %%%%%%%%%%%%%%%%%%%%%%%%%%%%%%%%%%%%%%%%%%%%%%%%%%%%%%%%%%%%%%%%%%%%%%%%%%%%%%%%%
%% Template for Research proposals for Computer Vision subject
%% %%%%%%%%%%%%%%%%%%%%%%%%%%%%%%%%%%%%%%%%%%%%%%%%%%%%%%%%%%%%%%%%%%%%%%%%%%%%%%%%%
%% Version 1.0 // February 2025
%% %%%%%%%%%%%%%%%%%%%%%%%%%%%%%%%%%%%%%%%%%%%%%%%%%
%% Eduardo FIDALGO
%%  - eduardo.fidalgo@unileon.es
%% %%%%%%%%%%%%%%%%%%%%%%%%%%%%%%%%%%%%%%%%%%%%%%%%%
\documentclass[11pt]{article}
\usepackage{AV_style}
\usepackage{comment}

\usepackage{lipsum}  %% Package to create dummy text (comment or erase before start)

\usepackage{multirow}
\usepackage{diagbox}
\usepackage{xfrac}
\usepackage{tikz}
\usepackage{caption, subcaption}
\usepackage{lscape}
\usepackage{pgfplots}
\usepackage[export]{adjustbox}
\usepackage{xcolor}

%% ===============================================
%% Setting the line spacing (3 options: only pick one)
% \doublespacing
% \singlespacing
\onehalfspacing
%% ===============================================

\setlength{\droptitle}{-5em} %% Don't touch

% %%%%%%%%%%%%%%%%%%%%%%%%%%%%%%%%%%%%%%%%%%%%%%%%%%%%%%%%%%
% SET THE TITLE
% %%%%%%%%%%%%%%%%%%%%%%%%%%%%%%%%%%%%%%%%%%%%%%%%%%%%%%%%%%
% TITLE:
\title{Research Proposal: Object detection and recognition of medical images 
}

% AUTHORS:
\author{Pablo Ruíz Morán\\% Student data
    \href{mailto:student.email@estudiantes.unileon.es}{\texttt{pruizm01@estudiantes.unileon.es}} %% Email author 1 
    }
    
% DATE:
%\date{\today}

% %%%%%%%%%%%%%%%%%%%%%%%%%%%%%%%%%%%%%%%%%%%%%%%%%%%%%%%%%%
% %%%%%%%%%%%%%%%%%%%%%%%%%%%%%%%%%%%%%%%%%%%%%%%%%%%%%%%%%%
\begin{document}
% %%%%%%%%%%%%%%%%%%%%%%%%%%%%%%%%%%%%%%%%%%%%%%%%%%%%%%%%%%
% %%%%%%%%%%%%%%%%%%%%%%%%%%%%%%%%%%%%%%%%%%%%%%%%%%%%%%%%%%
% ABSTRACT
% %%%%%%%%%%%%%%%%%%%%%%%%%%%%%%%%%%%%%%%%%%%%%%%%%%%%%%%%%%
% %%%%%%%%%%%%%%%%%%%%%%%%%%%%%%%%%%%%%%%%%%%%%%%%%%%%%%%%%%
{\setstretch{.8}
\maketitle
% %%%%%%%%%%%%%%%%%%
\begin{abstract}
This research project addresses the application of computer vision techniques to 
industrial inspection, focusing on the calibration and segmentation of cutting edges 
in rubber strip assembly within tyre manufacturing. The objective is to accurately 
detect and segment cutting borders, as well as estimate their spatial position and 
distances in a real-world industrial environment. This work is developed in the 
context of the \textbf{Michelin Challenge 2} proposed by \textit{Michelin Aranda de Duero}, where precision, robustness, and 
automation are critical for ensuring product quality and reducing operational costs.

To tackle this challenge, we analyze ten recent state-of-the-art methods (2024--2025) 
in object detection, image segmentation, and anomaly detection, with a particular 
focus on deep learning and vision-language models. Among the evaluated approaches, 
the Grounded-SAM framework is selected due to its ability to combine open-vocabulary 
grounding with high-quality segmentation. This method integrates Grounding DINO for 
text-guided object detection and the Segment Anything Model (SAM) for precise mask 
generation, enabling flexible detection of relevant regions without requiring 
task-specific training.

The selected deep learning method has demonstrated strong performance on 
open-vocabulary benchmarks, particularly on the Segmentation in the Wild (SegInW) 
zero-shot benchmark, which includes 25 diverse real-world datasets. Grounded-SAM 
achieves a mean Average Precision (mAP) of 48.7 in zero-shot settings, outperforming 
previous open-set segmentation approaches. These results highlight its strong 
generalization capabilities across heterogeneous and unseen scenarios, making it 
especially suitable for industrial environments characterized by variability, domain 
shifts, and limited annotated data.

Overall, this project aims to adapt and evaluate advanced vision-language models for 
cutting edge detection in a real industrial setting, incorporating calibration 
procedures to enable precise geometric measurements. The proposed approach 
contributes to the development of scalable, robust, and automated inspection systems 
for manufacturing processes.\\

\textit{\textbf{Keywords: } Cut Edge Detection, Object Detection, Image Segmentation, 
Grounded-SAM, Industrial Inspection, Michelin Challenge}
\end{abstract}
}

% %%%%%%%%%%%%%%%%%%%%%%%%%%%%%%%%%%%%%%%%%%%%%%%%%%%%%%%%%%
% %%%%%%%%%%%%%%%%%%%%%%%%%%%%%%%%%%%%%%%%%%%%%%%%%%%%%%%%%%
% BODY OF THE DOCUMENT
% %%%%%%%%%%%%%%%%%%%%%%%%%%%%%%%%%%%%%%%%%%%%%%%%%%%%%%%%%%
% %%%%%%%%%%%%%%%%%%%%%%%%%%%%%%%%%%%%%%%%%%%%%%%%%%%%%%%%%%

% --------------------
\section{Introduction}
\label{sec:introduction}
% Insert the content of the introduction here.
%    What is the problem?
%    Why is the problem significant or important?
%    A brief review of the state of the art.
%    What are the benefits or importance of this research?
%    What will be presented in this document?

Automatic detection and precise localization of cutting edges in industrial images is a key problem 
in modern manufacturing, driven by the increasing demand for automated 
quality control systems. In tyre assembly processes, rubber strips are placed over a cylindrical 
drum, cut, and subsequently joined at their edges. Currently, the tasks of verifying the position 
of the strip and joining its borders are performed manually by workers on the production floor, 
introducing variability, human error, and limiting the scalability of the process.

Precise cut edge localization is a particularly challenging problem in this industrial setting. 
Images are captured using smartphones by different workers from variable positions and angles, 
requiring the system to be robust to perspective distortion, changes in illumination, and 
variations in camera distance. Furthermore, the rubber strip and its cut edges often present low 
contrast with respect to the background, making accurate boundary detection non-trivial. An 
additional challenge is the requirement to provide metric measurements in millimetres rather 
than pixel coordinates, which necessitates a reliable pixel-to-metric calibration procedure based 
on QR codes placed on the inspection surface.

These challenges have been addressed through different paradigms in computer vision. Recent 
advances in deep learning have produced powerful methods for object detection and image 
segmentation. Convolutional approaches such as YOLOv8 \cite{Jocher2023YOLOv8} offer 
real-time detection with strong accuracy, while transformer-based detectors such as 
Relation-DETR \cite{HouRelationDETR2024} and RT-DETRv4 \cite{LiaoRTDETRv4} achieve 
state-of-the-art results on standard benchmarks. More recently, vision-language models such 
as YOLO-World \cite{ChengYOLOWorld2024} and Grounding DINO \cite{RenGroundedSAM2024} 
have extended object detection to open-vocabulary settings, enabling the detection of objects 
described by free-form text prompts without task-specific retraining. On the segmentation side, 
efficient transformer-based architectures such as Segformer++ \cite{KienzleSegformer2024} 
and PEM \cite{CavagneroPEM2024} provide accurate dense predictions, while the Segment 
Anything Model (SAM) \cite{KirillovSAM2023} enables promptable pixel-level segmentation 
for arbitrary objects. However, most of these methods are designed and evaluated on 
general-purpose benchmarks, and their direct applicability to industrial inspection tasks with 
limited labeled data remains an open challenge.

Despite this progress, applying these approaches to the specific problem of cut edge detection 
in tyre manufacturing introduces additional difficulties, including the absence of large annotated 
datasets, the need for reliable geometric calibration under uncontrolled acquisition conditions, 
and the requirement for real-time or near-real-time performance in production environments.

The main contribution of this work is the design and adaptation of open-vocabulary vision-language models 
to a real industrial inspection scenario, integrating segmentation and geometric calibration into 
a unified framework. In particular, we analyze ten recent state-of-the-art methods for object 
detection, segmentation, and anomaly detection, and select Grounded-SAM 
\cite{RenGroundedSAM2024} as the primary deep learning approach. Grounded-SAM combines 
Grounding DINO for open-vocabulary detection with SAM for precise mask generation, enabling 
flexible detection and segmentation of cut edges without requiring task-specific training data. 
Additionally, the proposed system incorporates QR-based spatial calibration to correct 
perspective distortions and provide metric distance estimations, contributing to the development 
of a scalable and automated inspection pipeline for the Michelin production environment.

The rest of the document is organized as follows. Section~\ref{sec:literatureReview} presents 
a review of the ten selected state-of-the-art methods. Section~\ref{sec:selectedMethod} 
describes the selected methods and the full computer vision pipeline proposed for 
the challenge. Section~\ref{sec:datasets} presents the datasets used for evaluation, and 
Section~\ref{sec:metrics} defines the evaluation metrics. Section~\ref{sec:python} provides 
the Python implementations available for the reviewed methods. Finally, 
Section~\ref{sec:credits} lists the individual contributions of each group member.
% --------------------
% --------------------
\section{Literature Review}
% --------------------
\label{sec:literatureReview}
% Insert the content of the proposed solutions here.
% Papers reviewed must be recent (2024-2025), no more than three years old.

The following section reviews ten recent methods (2024--2025) relevant to the \textbf{Michelin Challenge}.
The reviewed methods can be broadly categorized into Object Detection, Image Segmentation and Anomaly Detection.\\

\textit{Jocher et al. \cite{Jocher2023YOLOv8} presents an in-depth analysis of YOLOv8, a 
next-generation object detection model developed by Ultralytics. The architecture improves 
previous YOLO versions by integrating a CSPNet backbone for enhanced feature extraction, an 
FPN+PAN neck for multi-scale feature aggregation, and an anchor-free detection head that 
simplifies bounding box prediction and reduces computational complexity. The model also 
introduces improved training strategies such as advanced data augmentation techniques (mosaic 
and mixup), focal loss for classification, and mixed-precision training to accelerate learning on 
modern GPUs. Benchmark evaluations on Microsoft COCO and Roboflow 100 demonstrate that 
YOLOv8 achieves higher accuracy and faster inference compared with earlier models like 
YOLOv5, reaching $71.5\%$ mAP@0.5 with the YOLOv8x variant at 11.5 ms inference time.}\\

\textit{Liao et al.~\cite{LiaoRTDETRv4} proposed RT-DETRv4, a real-time object detection 
framework that further improves DETR-based detectors by distilling semantic knowledge from 
Vision Foundation Models into lightweight detection architectures. The method introduces a 
Deep Semantic Injector (DSI) module to inject high-level semantic representations into the 
deep layers of the detector, together with a Gradient-guided Adaptive Modulation (GAM) 
strategy that dynamically balances semantic transfer during training. Evaluated on COCO 2017, 
RT-DETRv4 achieved AP scores of 49.7, 53.5, 55.4, and 57.0 for the S, M, L, and X variants 
respectively, at corresponding speeds of 273, 169, 124, and 78 FPS.}\\

\textit{Wyatt et al.~\cite{WyattAnoDDPM2022} proposed an unsupervised anomaly detection 
approach for defect localization in ultrasonic non-destructive testing using a generative 
diffusion model. Their method relies on AnoDDPM trained exclusively on defect-free data to 
reconstruct normal ultrasonic wave propagation patterns. During inference, defects are detected 
by computing the difference between the input image and the reconstructed output, producing 
anomaly maps that highlight defective regions. Applied to a Laser Ultrasonic Visualization 
Testing (LUVT) dataset for aluminium inspection, the method achieved an AUROC of 0.909, 
outperforming VAE ($0.800$) and f-AnoGAN ($0.755$), and reaching Precision $0.817$ and 
Recall $0.820$ when compared against supervised detectors such as YOLOv3 and RetinaNet.}\\

\textit{Kienzle et al.~\cite{KienzleSegformer2024} proposed Segformer++, an efficient 
transformer-based architecture for high-resolution semantic segmentation that adapts 
token-merging strategies to the Segformer framework without requiring model retraining. The 
authors introduce two configurations: Segformer++HQ, oriented towards quality preservation, 
and Segformer++fast, oriented towards speed. Evaluated on Cityscapes and ADE20K, 
Segformer++HQ achieves a mIoU of $82.31$ on Cityscapes with a $1.61\times$ inference 
speedup, while Segformer++fast reaches a $1.94\times$ speedup with only a marginal 
performance drop.}\\

\textit{Cavagnero et al.~\cite{CavagneroPEM2024} proposed PEM (Prototype-based Efficient 
MaskFormer), an efficient transformer-based architecture for image segmentation that 
introduces a prototype-based cross-attention mechanism to reduce computational burden, 
together with an efficient multi-scale feature pyramid network combining deformable 
convolutions and context-based self-modulation. PEM was evaluated on both semantic and 
panoptic segmentation tasks using Cityscapes and ADE20K, achieving $79.9$ mIoU on 
Cityscapes at 13.8 FPS for semantic segmentation and $61.1$ PQ on Cityscapes at 13.8 FPS 
for panoptic segmentation, demonstrating a strong trade-off between segmentation accuracy 
and efficiency.}\\

\textit{Hou et al.~\cite{HouRelationDETR2024} proposed Relation-DETR, a transformer-based 
object detector that introduces an explicit position relation prior to improve both convergence 
speed and detection accuracy in DETR architectures. The method incorporates a position 
relation encoder that models pairwise interactions between bounding boxes and uses this 
information as an attention bias for progressive attention refinement. In addition, the authors 
extend the standard DETR pipeline into a contrast relation pipeline combining one-to-one and 
one-to-many matching. Evaluated on COCO 2017, Relation-DETR achieved $51.7$ AP with 
ResNet-50 in the $1\times$ training setting, outperforming DINO by $+2.0$ AP and reaching 
over 40 AP after only 2 training epochs.}\\

\textit{Heidari et al.~\cite{HeidariViLUNet2025} introduced ViLU-Net, a hybrid 
encoder--decoder architecture that integrates Vision-LSTM (xLSTM) blocks within the 
traditional U-Net structure to model long-range spatial dependencies while maintaining 
efficient computational complexity. The model combines CNN layers for local feature extraction 
with xLSTM modules for global contextual information. Evaluated on an in-house CT dataset 
of retroperitoneal tumours and the FLARE abdominal organ segmentation dataset, ViLU-Net 
achieves a Dice Similarity Coefficient (DSC) of $0.9309$ and IoU of $0.8720$, outperforming 
methods such as nnU-Net, SwinUNETR, and U-Mamba.}\\

\textit{Ren et al.~\cite{RenGroundedSAM2024} introduced Grounded-SAM, a modular framework 
that combines Grounding DINO, an open-set object detector capable of identifying objects based 
on free-form text prompts, with the Segment Anything Model (SAM), a promptable segmentation 
model that generates precise pixel-level masks. In the proposed pipeline, Grounding DINO 
detects candidate regions and outputs bounding boxes, which are then used as prompts by SAM 
to generate accurate segmentation masks. Evaluated on the Segmentation in the Wild (SegInW) 
benchmark --- comprising 25 diverse real-world datasets --- the method achieves a mean Average 
Precision of $48.7$ in zero-shot settings, demonstrating strong generalization capabilities for 
open-world vision tasks without additional training.}\\

\textit{Cheng et al.~\cite{ChengYOLOWorld2024} propose a real-time open-vocabulary 
object detector that enables the identification of objects not explicitly 
included in the training set. The model extends the YOLO architecture by 
incorporating linguistic information through a CLIP-based text encoder and 
a Re-parameterizable Vision-Language Path Aggregation Network (RepVL-PAN), 
which fuses visual and semantic features. As a result, YOLO-World achieves 
strong performance in open-world scenarios, reaching $35.4$ AP on the LVIS 
dataset in zero-shot settings at 52 FPS.}\\

\textit{Ren et al.~\cite{RenGroundingDINO15} introduced Grounding DINO 1.5, an advanced 
open-vocabulary object detection model that extends the original Grounding DINO framework. 
It combines visual features with textual prompts through a transformer-based architecture, 
enabling detection of objects beyond predefined categories. Trained on large-scale grounding 
datasets, the Pro model achieves $54.3$ AP on COCO and $55.7$ AP on LVIS-minival in 
zero-shot settings, while an optimised Edge variant reaches 75.2 FPS on A100 with TensorRT, 
making it suitable for industrial deployment scenarios.}\\

A summary of the reviewed methods is provided in Table~\ref{tab:literatureReview}.

\begin{table}[phtb!]
\begin{center}
\resizebox{\linewidth}{!}{
\begin{tabular}{l c p{4.5cm} c c c c}
\hline
\multirow{2}{*}{Method} & \multirow{2}{*}{Year} & \multirow{2}{*}{Description} & 
\multirow{2}{*}{Dataset} & \multirow{2}{*}{Main Metric} & 
\multirow{2}{*}{Code} & \multirow{2}{*}{Speed} \\
 & & & & & & \\
\hline
YOLOv8 \cite{Jocher2023YOLOv8} & 2024 & 
    Anchor-free CNN detector with CSPNet backbone and FPN+PAN neck & 
    COCO, Roboflow 100 & 71.5\% mAP@0.5 & Yes & 11.5 ms \\
RT-DETRv4 \cite{LiaoRTDETRv4} & 2025 & 
    DETR detector with VFM distillation via DSI and GAM & 
    COCO 2017 & 49.7--57.0 AP & Yes & 78--273 FPS \\
AnoDDPM \cite{WyattAnoDDPM2022} & 2024 & 
    Unsupervised anomaly detection with diffusion models & 
    LUVT & AUROC 0.909 & Yes & N/R \\
Segformer++ \cite{KienzleSegformer2024} & 2024 & 
    Efficient transformer-based semantic segmentation via token merging & 
    Cityscapes, ADE20K & 82.31 mIoU & Yes & $1.61\times$ speedup \\
PEM \cite{CavagneroPEM2024} & 2024 & 
    Prototype-based efficient MaskFormer & 
    Cityscapes, ADE20K & 61.1 PQ / 79.9 mIoU & Yes & 13.8--35.7 FPS \\
Relation-DETR \cite{HouRelationDETR2024} & 2024 & 
    DETR with explicit position relation prior & 
    COCO 2017 & 51.5 AP & Yes & N/R \\
ViLU-Net \cite{HeidariViLUNet2025} & 2025 & 
    Hybrid U-Net with CNN and xLSTM blocks & 
    CT Tumour, FLARE & DSC 0.9309 / IoU 0.8720& Yes & N/R \\
Grounded-SAM \cite{RenGroundedSAM2024} & 2024 & 
    Open-vocabulary detection + segmentation pipeline & 
    SegInW (zero-shot) & 48.7 mAP & Yes & N/R \\
YOLO-World \cite{ChengYOLOWorld2024} & 2024 & 
    Real-time open-vocabulary YOLO with RepVL-PAN & 
    LVIS & 28.7 AP / 35.2 APf & Yes & 52.0 FPS \\
Grounding DINO 1.5 \cite{RenGroundingDINO15} & 2024 & 
    Advanced open-vocabulary detector with ViT-L backbone & 
    COCO, LVIS & 54.3 AP, 55.7 AP & Yes & 5.5--75.2 FPS \\
\hline
\end{tabular}
}
\caption{\label{tab:literatureReview}Summary of the reviewed methods for object detection, 
segmentation, and anomaly detection. N/R = Not Reported.}
\end{center}
\end{table}

Overall, the reviewed methods show that transformer-based approaches achieve 
state-of-the-art performance in both object detection and segmentation tasks, while 
diffusion-based methods provide promising results for unsupervised anomaly detection. 
However, transformer models often require high computational resources, whereas 
YOLO-based approaches offer a better trade-off between speed and accuracy for 
real-time applications. Open-vocabulary models such as Grounded-SAM and YOLO-World 
demonstrate strong generalization capabilities, making them suitable for real-world 
industrial scenarios where object variability and limited labeled data are key constarints.\\

%%%%%%%%%%%%%%%%%%%%%%%%%%%%%%%%%%%%%%%
%%%%%%%%%%%%%%%%%%%%%%%%%%%%%%%%%%%%%%%
% METODOS JAIME
%%%%%%%%%%%%%%%%%%%%%%%%%%%%%%%%%%%%%%%
%%%%%%%%%%%%%%%%%%%%%%%%%%%%%%%%%%%%%%%

Ravi et al. \cite{Ravi2024SAM2} presented SAM 2 (Segment Anything Model 2), 
a foundation model developed by Meta FAIR for promptable visual segmentation 
in both images and videos. The architecture consists of a hierarchical Hiera 
image encoder pretrained with MAE, a prompt encoder that accepts points, 
bounding boxes or masks, a two-way transformer mask decoder, and a streaming 
memory module composed of a memory encoder, a FIFO memory bank and a memory 
attention block that enables segmentation propagation across frames. Trained 
on the SA-V dataset\footnote{https://ai.meta.com/datasets/segment-anything-video/}, 
comprising 50.9K videos and 35.5M masks, SAM 2 achieves a J\&F of $76.6\%$ 
on MOSE val using the Hiera-B+ encoder, while being $6\times$ faster than 
its predecessor SAM on image segmentation benchmarks, obtaining $61.9\%$ 
mIoU across 23 zero-shot datasets.\\

Xiong et al. \cite{Xiong2023EfficientSAM} proposed EfficientSAM, a lightweight 
version of SAM developed by Meta AI Research that addresses the high computational 
cost of the original model. Their approach, SAMI (SAM-leveraged Masked Image 
Pretraining), trains lightweight ViT encoders (ViT-Tiny and ViT-Small) to 
reconstruct feature embeddings from SAM's ViT-H image encoder using a 
cross-attention decoder and MSE reconstruction loss, transferring SAM's 
representational knowledge into compact models. The resulting EfficientSAM-S 
reduces SAM's parameter count by $\sim20\times$ (25M vs 636M) while achieving 
$44.4$ AP on zero-shot instance segmentation on COCO\footnote{https://cocodataset.org}, 
compared to SAM's $46.5$ AP, at a throughput of 47 images/second versus SAM's 2 images/second.\\

Wang et al. \cite{Wang2024YOLOv9} proposed YOLOv9, a real-time object detection 
model that addresses the information bottleneck problem in deep networks through 
two key contributions. First, Programmable Gradient Information (PGI), an 
auxiliary supervision framework that uses a reversible branch to generate reliable 
gradients and preserve complete input information during training without adding 
inference cost. Second, the Generalized Efficient Layer Aggregation Network 
(GELAN), a lightweight architecture based on CSPNet and ELAN that supports 
arbitrary computational blocks while maintaining strong parameter efficiency. 
Evaluated on MS COCO\footnote{\url{https://cocodataset.org}}, YOLOv9-E achieves 
$55.6\%$ AP$_{50:95}$ with 57.3M parameters, reducing parameter count by $49\%$ 
and computation by $43\%$ compared to YOLOv8-X while improving AP by $0.6\%$.\\

Shi et al. \cite{Shi2025CarbonBlock} proposed YOLOv8-HDSA, an improved 
YOLOv8-based instance segmentation algorithm for industrial carbon block 
recognition developed at Shandong University of Science and Technology. 
The architecture introduces three main contributions over the baseline 
YOLOv8: a Selective Reinforcement Feature Fusion Module (SRFF) that uses 
Hadamard product and dilated convolution to enhance feature representation 
while suppressing background noise; a convolutional self-attention mechanism 
with residual structure added to the detection head to improve fine-grained 
edge detail extraction; and Focaler-IoU as a loss function to dynamically 
adjust sample weights and optimize regression performance. Evaluated on a 
real industrial carbon block dataset\footnote{\url{https://www.nature.com/articles/s41598-025-91495-x}}, YOLOv8-HDSA improved average recognition accuracy by $7.2\%$ and segmentation 
accuracy by $3.8\%$ over the baseline YOLOv8.\\

Damm et al. \cite{Damm2024AnomalyDINO} proposed AnomalyDINO, a training-free 
vision-only framework for one-shot and few-shot anomaly detection in industrial 
applications. The method adapts DINOv2 as a patch-level feature extractor and 
follows the deep nearest neighbor paradigm, computing distances between patch 
embeddings of test images and a small set of normal reference images to produce 
both image-level anomaly scores and pixel-level anomaly segmentation maps, 
without requiring fine-tuning, meta-learning, or any additional training data. 
Evaluated on the MVTec-AD dataset\footnote{\url{https://www.mvtec.com/company/research/datasets/mvtec-ad}}, 
AnomalyDINO achieves an AUROC of $96.6\%$ in the one-shot setting, improving 
over the previous state-of-the-art by $3.5\%$, while also outperforming 
existing vision-language approaches in multiple few-shot configurations.\\

Jiang et al. \cite{Jiang2025CLIPDINO} proposed a zero-shot industrial anomaly 
detection framework that fuses CLIP's global semantic embeddings with DINOv2's 
multi-scale structural features through a Dual-Modality Attention (DMA) mechanism. 
The DMA module computes DINOv2 self-similarity weights and CLIP-guided attention 
maps that are adaptively fused to align cross-modal representations without 
requiring any target-domain training samples. A complementary Stabilized 
Attention-based Pooling (SAP) module uses self-generated anomaly heatmaps as 
prior guidance to aggregate discriminative features and suppress normal background 
dilution. Trained with a multi-task loss combining Focal, Dice, and contrastive 
losses on auxiliary datasets, the framework is evaluated across seven industrial 
benchmarks\footnote{\url{https://www.mvtec.com/company/research/datasets/mvtec-ad}}, 
achieving an average image-level AUROC of $93.4\%$, AP of $94.3\%$, 
pixel-level AUROC of $96.9\%$, and AUPRO of $92.4\%$, outperforming 
state-of-the-art zero-shot methods including WinCLIP, AnomalyCLIP, and Crane.\\

Xiong et al. \cite{Xiong2026SAM2UNet} proposed SAM2-UNet, a simple yet 
effective U-shaped segmentation framework that leverages the hierarchical 
Hiera backbone of SAM2 as encoder. Unlike the plain ViT encoder used in 
SAM1, Hiera provides multi-scale features at four resolution levels, making 
it naturally suited for a U-Net decoder architecture. Lightweight adapters 
with a bottleneck MLP design are inserted before each Hiera block to enable 
parameter-efficient fine-tuning while keeping the encoder frozen. The decoder 
follows the classic U-Net design with three decoder blocks and deep supervision, 
trained with weighted IoU and binary cross-entropy losses. Evaluated across 
18 datasets\footnote{\url{https://github.com/WZH0120/SAM2-UNet}} on five 
benchmarks including camouflaged object detection, salient object detection, 
marine animal segmentation, mirror detection, and polyp segmentation, 
SAM2-UNet consistently outperforms specialized state-of-the-art methods, 
achieving for example a mDice of $92.8\%$ on Kvasir-SEG and an S-measure 
improvement of $4.8\%$ over FEDER on the CAMO dataset.\\

Shi et al. \cite{Shi2024SemiSupervised} proposed a semi-supervised semantic 
segmentation framework for industrial surface defect detection that addresses 
the scarcity of labeled data through a Two-Branch Perturbation Cross Pseudo 
Supervision model (TPcps). One branch applies image-level perturbations 
(geometric and non-geometric augmentations, noise injection, and random 
erasure), while the second branch perturbs the feature space via random 
channel dropout, with both branches generating mutual pseudo-labels constrained 
by a cross pseudo-supervised loss. The segmentation network combines a 
lightweight modified MobileNetv2 backbone with DeepLabv3Plus, augmented by 
a lightweight attention module using 1D convolution for channel attention 
and dilated convolution for spatial attention, plus a shallow feature fusion 
module. A dynamic threshold strategy prevents simultaneous false label 
generation between branches. Evaluated on five self-built industrial 
datasets\footnote{\url{https://doi.org/10.1038/s41598-024-72579-6}} and the 
public KolektorSDD dataset, the method achieves a mIoU $3.11\%$ higher than 
the state-of-the-art on the self-built dataset and $4.39\%$ higher on 
KolektorSDD, while matching the inference speed of U-Net at $64.59$ fps.\\

Liao et al. \cite{Liao2025CameraCalibration} presented a comprehensive survey 
of learning-based camera calibration techniques covering a decade of research. 
The work categorizes existing methods into four main models: the standard 
pinhole camera model, the distortion camera model, the cross-view model, and 
the cross-sensor model, analyzing deep learning strategies including 
convolutional networks, transformers, and geometric prior integration. The 
survey reviews approaches for intrinsic and extrinsic parameter estimation, 
lens distortion correction, homography estimation, and multi-camera 
calibration, highlighting how learning-based methods overcome the limitations 
of traditional checkerboard-based pipelines. To address the lack of a unified 
benchmark, the authors compile a holistic calibration 
dataset\footnote{\url{https://github.com/KangLiao929/Awesome-Deep-Camera-Calibration}} 
comprising both synthetic and real-world images captured by diverse cameras 
in varied scenes. This survey is directly applicable to the Michelin challenge, 
where QR code-based distance calibration from mobile phone images with variable 
camera position requires robust geometric parameter estimation under 
perspective distortion.\\

Heckler-Kram et al. \cite{HecklerKram2025MVTecAD2} introduced MVTec AD 2, 
a new benchmark dataset for unsupervised anomaly detection in industrial 
visual inspection that addresses the saturation of existing benchmarks such 
as MVTec AD and VisA, where state-of-the-art models compete within less than 
one percentage point of AU-PRO. The dataset comprises eight anomaly detection 
scenarios with over 8000 high-resolution 
images\footnote{\url{https://www.mvtec.com/company/research/datasets/mvtec-ad-2}}, 
covering challenging industrial use cases not addressed by prior datasets, 
including transparent and overlapping objects, dark-field and backlight 
illumination, objects with high variance in normal data, and extremely small 
defects distributed across the full image including borders. Each scenario 
is captured under multiple lighting conditions to evaluate model robustness 
to real-world distribution shifts. Comprehensive evaluation of state-of-the-art 
methods shows that performance remains below $60\%$ average AU-PRO, 
demonstrating the dataset's discriminatory power for benchmarking future 
industrial anomaly detection methods.\\

\begin{table}[phtb!]
\begin{center}
\resizebox{\linewidth}{!}{
\begin{tabular}{l c p{4.5cm} c c c c}
\hline
\multirow{2}{*}{Method} & \multirow{2}{*}{Year} & \multirow{2}{*}{Description} & 
\multirow{2}{*}{Dataset} & \multirow{2}{*}{Main Metric} & 
\multirow{2}{*}{Code} & \multirow{2}{*}{Speed} \\
 & & & & & & \\
\hline
SAM 2 \cite{Ravi2024SAM2} & 2024 & 
    Transformer with streaming memory for promptable image and video segmentation & 
    MOSE val, SA-23 & J\&F 76.6\% / mIoU 61.9\% & Yes & 43.8 FPS \\
EfficientSAM \cite{Xiong2023EfficientSAM} & 2023 & 
    Lightweight SAM via SAMI pretraining with ViT-Tiny/Small encoders & 
    COCO, LVIS & 44.4 AP (zero-shot) & Yes & 47--54 img/s \\
YOLOv9 \cite{Wang2024YOLOv9} & 2024 & 
    Real-time detector with PGI and GELAN for gradient information preservation & 
    MS COCO & 55.6\% AP$_{50:95}$ & Yes & N/R \\
YOLOv8-HDSA \cite{Shi2025CarbonBlock} & 2025 & 
    Improved YOLOv8 with SRFF module and Focaler-IoU for industrial segmentation & 
    Industrial carbon block & +7.2\% accuracy over YOLOv8 & Yes & N/R \\
AnomalyDINO \cite{Damm2024AnomalyDINO} & 2024 & 
    Training-free few-shot anomaly detection via DINOv2 patch nearest neighbor & 
    MVTec-AD & AUROC 96.6\% (1-shot) & Yes & N/R \\
CLIP-DINOv2 \cite{Jiang2025CLIPDINO} & 2025 & 
    Zero-shot anomaly detection with DMA fusion of CLIP and DINOv2 features & 
    MVTec AD, VisA (7 benchmarks) & AUROC 93.4\% / AP 94.3\% & No & 23 FPS \\
SAM2-UNet \cite{Xiong2026SAM2UNet} & 2026 & 
    U-shaped segmentation using SAM2 Hiera encoder with PEFT adapters & 
    18 datasets (5 benchmarks) & mDice 92.8\% (Kvasir) & Yes & N/R \\
Semi-sup. defect seg. \cite{Shi2024SemiSupervised} & 2024 & 
    Two-branch cross pseudo-supervision with lightweight DeepLabv3Plus for defects & 
    KolektorSDD, Steel & mIoU +4.39\% over SOTA & No & 64.59 FPS \\
Camera Calib. Survey \cite{Liao2025CameraCalibration} & 2025 & 
    Survey of DL-based camera calibration: pinhole, distortion, cross-view models & 
    Holistic calib. dataset & N/R (survey) & Yes & N/R \\
MVTec AD 2 \cite{HecklerKram2025MVTecAD2} & 2025 & 
    Advanced industrial anomaly detection benchmark with 8 challenging scenarios & 
    MVTec AD 2 & AU-PRO $<$60\% (SOTA) & No & N/R \\
\hline
\end{tabular}
}
\caption{\label{tab:literatureReview}Summary of the reviewed methods for 
segmentation, anomaly detection, and camera calibration. N/R = Not Reported.}
\end{center}
\end{table}

% --------------------
% --------------------
\section{Selected Method}
% --------------------
\label{sec:selectedMethod}
% Each student must describe one method in detail (at least one traditional 
% and one deep learning approach per group).
% TODO: add subsections for the traditional method(s) selected by other group members.

Grounded-SAM \cite{RenGroundedSAM2024} has been selected as the primary deep 
learning due to its ability to combine 
open-vocabulary object detection with high-quality segmentation in a unified pipeline.

% -----------------------------------------------
\subsection{Grounded-SAM (Deep Learning Approach)}
% -----------------------------------------------

\subsubsection{Description and motivation}

Grounded-SAM \cite{RenGroundedSAM2024} is a computer vision framework that combines 
Grounding DINO, an open-vocabulary object detector, with the Segment Anything Model (SAM) 
to perform automatic object detection and segmentation. The method enables detecting objects 
using natural language prompts and generating precise pixel-level segmentation masks, without 
requiring task-specific retraining.

This method was selected for the Michelin Challenge 2 because the combination of detection 
and segmentation in a single pipeline allows the system to:
\begin{itemize}
    \item detect QR codes for spatial calibration,
    \item segment the table or conveyor belt area,
    \item detect and segment rubber strip cut edges,
    \item estimate distances between detected elements using geometric analysis.
\end{itemize}

\subsubsection{Architecture}

Grounded-SAM has a modular design composed of two main deep learning models working sequentially.

\paragraph{Grounding DINO (Object Detection Stage).}
Grounding DINO is a transformer-based object detector that performs open-vocabulary 
detection using natural language prompts. The transformer aligns visual features with text 
embeddings through cross-attention mechanisms, allowing the model to associate image 
regions with semantic concepts. Its main components are:

\begin{itemize}
    \item \textbf{Backbone network:} Extracts multi-scale visual features from the input image.
    \item \textbf{Text encoder:} Encodes the input text prompt (e.g., ``QR code'' or 
    ``rubber strip edge'') into text embeddings.
    \item \textbf{Transformer encoder-decoder:} Processes visual features and aligns them 
    with text embeddings through cross-attention layers.
    \item \textbf{Detection head:} Predicts bounding boxes and confidence scores 
    corresponding to the detected objects.
\end{itemize}

The output of this stage is a set of bounding boxes associated with the objects described 
in the text prompt.

\paragraph{Segment Anything Model (SAM) (Segmentation Stage).}
SAM performs high-quality segmentation based on input prompts. Its architecture includes:

\begin{itemize}
    \item \textbf{Image encoder (Vision Transformer):} Extracts high-level visual feature 
    embeddings from the entire input image. This stage runs only once per image.
    \item \textbf{Prompt encoder:} Encodes prompts such as bounding boxes or points 
    into dense and sparse embeddings.
    \item \textbf{Mask decoder:} Generates the final pixel-level segmentation mask by 
    combining image embeddings with prompt embeddings through cross-attention and 
    upsampling operations.
\end{itemize}

SAM receives the bounding boxes produced by Grounding DINO as spatial prompts and 
produces pixel-level segmentation masks for the detected objects.

\subsubsection{Computer vision pipeline}

\paragraph{Step 1: Image acquisition.}
An RGB image of the inspection area is captured using a smartphone. Since images are 
taken by different workers at variable positions and angles, the system must be robust to 
perspective distortion and changes in illumination.

\paragraph{Step 2: Preprocessing.}
Before feeding the image into the neural networks, the following operations are applied:
\begin{itemize}
    \item \textbf{Image resizing:} The image is resized to match the input resolution 
    expected by each model.
    \item \textbf{Pixel normalization:} Pixel values are normalized to follow the 
    distribution used during training.
    \item \textbf{Tensor conversion:} The image is converted into a PyTorch tensor for 
    neural network processing.
    \item \textbf{QR code detection:} QR codes in the scene are detected and decoded 
    to extract the calibration reference points for metric distance estimation.
\end{itemize}

\paragraph{Step 3: Feature extraction.}
The preprocessed image is processed by the Grounding DINO backbone, which extracts 
multi-scale visual feature maps capturing both global and local visual information from 
the scene.

\paragraph{Step 4: Object detection with Grounding DINO.}
Grounding DINO receives the visual features and a text prompt describing the objects 
of interest. Examples of prompts used in this application include:
\begin{itemize}
    \item \texttt{"QR code"}
    \item \texttt{"rubber strip"}
    \item \texttt{"cutting edge"}
\end{itemize}
The output of this stage is a set of bounding boxes localizing the objects described 
by the prompt.

\paragraph{Step 5: Segmentation with SAM.}
The bounding boxes produced by Grounding DINO are passed as spatial prompts to SAM, 
which generates pixel-level segmentation masks for the detected objects.

\paragraph{Step 6: Edge extraction and post-processing.}
Once segmentation masks are obtained, image processing techniques are applied to 
extract geometric information:
\begin{itemize}
    \item contour detection,
    \item edge extraction,
    \item boundary analysis.
\end{itemize}
These techniques allow identifying the exact borders of the segmented rubber strips 
or cutting edges.

\paragraph{Step 7: Geometric analysis and metric estimation.}
Using the segmented masks and the QR codes as reference markers, the system computes:
\begin{itemize}
    \item distances between cutting edges,
    \item distance from the edge to the table border,
    \item object positions in the coordinate system defined by the QR codes, where 
    point $(0, 0)$ is the centre of the upper-left QR code.
\end{itemize}
The pixel-to-millimetre conversion is derived from the known physical distance between 
QR codes (800 mm legs of a right triangle) using homography or affine transformations 
to correct for perspective distortion.

\subsubsection{Output}

The output of the Grounded-SAM pipeline includes:
\begin{itemize}
    \item bounding boxes of detected objects,
    \item pixel-level segmentation masks of the rubber strip cut edges,
    \item pixel coordinates of detected edge boundaries,
    \item metric measurements of edge positions and distances in millimetres.
\end{itemize}

\subsubsection{Evaluation metrics}

The performance of Grounded-SAM is evaluated using the following metrics. For the object 
detection stage performed by Grounding DINO: Average Precision (AP) and mean Average 
Precision (mAP), which evaluate the accuracy of predicted bounding boxes with respect to 
ground truth annotations. For the segmentation stage performed by SAM: Intersection over 
Union (IoU) and mask IoU, which measure the overlap between the predicted segmentation 
mask and the ground truth mask. These performance metrics are described in detail in
Section~\ref{sec:metrics}

\subsubsection{Applicability to the Michelin Challenge}

Grounded-SAM is particularly suited for the Michelin Challenge 2 because it integrates 
object detection and segmentation in a unified pipeline without requiring task-specific 
training data. In this context, the method can detect QR codes for spatial calibration, 
segment the table or conveyor belt area, identify the cutting edges of rubber strips, and 
compute metric distances between detected elements. The combination of open-vocabulary 
detection and precise segmentation enables accurate localization of relevant structures 
and supports the measurement tasks required in the automated inspection system.

\subsection{YOLOv9 + EfficientSAM (Deep Learning Approach)}

\subsubsection{Description and motivation}

YOLOv9 + EfficientSAM is a two-stage computer vision pipeline that combines a 
real-time object detector with a lightweight segmentation model to perform accurate 
detection and pixel-level segmentation of rubber strip cut edges. YOLOv9 
\cite{WangYOLOv92024} is a state-of-the-art real-time object detector that introduces 
Programmable Gradient Information (PGI) and the Generalized Efficient Layer Aggregation 
Network (GELAN) to address the information bottleneck problem in deep networks. 
EfficientSAM \cite{XiongEfficientSAM2023} is a lightweight variant of the Segment 
Anything Model (SAM) that leverages masked image pretraining to achieve high-quality 
segmentation with significantly reduced computational cost.

This method was selected for the Michelin Challenge 2 because the combination of 
detection and segmentation in a sequential pipeline allows the system to:
\begin{itemize}
    \item detect QR codes for spatial calibration and perspective correction,
    \item segment the table and rubber band area,
    \item detect and segment the cut edges of the rubber strip,
    \item estimate metric distances between detected elements using geometric analysis.
\end{itemize}

\subsubsection{Architecture}

The pipeline is composed of two main deep learning models working sequentially.

\paragraph{YOLOv9 (Object Detection Stage).}
YOLOv9 \cite{WangYOLOv92024} is a real-time object detector built on two core 
contributions:

\begin{itemize}
    \item \textbf{GELAN (Generalized Efficient Layer Aggregation Network):} A lightweight 
    network architecture based on gradient path planning that combines CSPNet and ELAN 
    blocks. GELAN achieves better parameter utilization than depth-wise convolution-based 
    designs while maintaining high inference speed and accuracy.
    \item \textbf{PGI (Programmable Gradient Information):} An auxiliary supervision 
    framework composed of an auxiliary reversible branch and multi-level auxiliary 
    information. PGI ensures that reliable gradient information is available throughout 
    training, preventing information loss caused by the information bottleneck in deep 
    networks. The auxiliary reversible branch is removed at inference time, so it 
    introduces no additional computational cost.
\end{itemize}

The main components of YOLOv9 are:
\begin{itemize}
    \item \textbf{Backbone (GELAN):} Extracts multi-scale visual features from the input 
    image while retaining complete information through gradient path planning.
    \item \textbf{Neck (PAN):} Aggregates features from different scales to improve 
    detection of objects of varying sizes.
    \item \textbf{Detection head:} Predicts bounding boxes and confidence scores for 
    the detected objects.
\end{itemize}

The output of this stage is a set of bounding boxes localizing the cut edges in 
the image.

\paragraph{EfficientSAM (Segmentation Stage).}
EfficientSAM \cite{XiongEfficientSAM2023} is a lightweight variant of SAM obtained 
through a pretraining method called SAMI (SAM-leveraged Masked Image pretraining). 
Instead of reconstructing raw pixel values as in standard masked autoencoders, SAMI 
trains a lightweight Vision Transformer encoder to reconstruct the feature embeddings 
produced by the SAM ViT-H image encoder. This allows the lightweight encoder to acquire 
the strong visual representations of SAM at a fraction of the computational cost. 
Its architecture includes:

\begin{itemize}
    \item \textbf{Lightweight image encoder (ViT-Tiny or ViT-Small):} Pretrained with 
    SAMI to extract high-quality visual feature embeddings from the input image.
    \item \textbf{Prompt encoder:} Encodes spatial prompts such as bounding boxes or 
    points into dense and sparse embeddings.
    \item \textbf{Mask decoder:} Generates the final pixel-level segmentation mask by 
    combining image embeddings with prompt embeddings through cross-attention and 
    upsampling operations.
\end{itemize}

EfficientSAM receives the bounding boxes produced by YOLOv9 as spatial prompts and 
generates pixel-level segmentation masks for the detected cut edges.

\subsubsection{Computer vision pipeline}

\paragraph{Step 1: Image acquisition.}
An RGB image of the inspection area is captured using a smartphone. Since images are 
taken by different workers at variable positions and angles, the system must be robust 
to perspective distortion and illumination changes.

\paragraph{Step 2: Preprocessing.}
Before feeding the image into the neural networks, the following operations are applied:
\begin{itemize}
    \item \textbf{QR code detection and decoding:} The three QR codes present in the 
    scene are detected and decoded using OpenCV. Their pixel coordinates are extracted 
    to define the spatial reference system, where point $(0, 0)$ corresponds to the 
    centre of the upper-left QR code.
    \item \textbf{Homography estimation:} Using the known physical distance between 
    QR codes (800 mm legs of a right triangle) and their detected pixel positions, a 
    homography matrix is computed to correct perspective distortion and derive the 
    pixel-to-millimetre conversion factor.
    \item \textbf{Image resizing and normalization:} The image is resized to match the 
    input resolution expected by each model, pixel values are normalized, and the image 
    is converted into a PyTorch tensor.
\end{itemize}

\paragraph{Step 3: Feature extraction.}
The preprocessed image is processed by the YOLOv9 GELAN backbone, which extracts 
multi-scale visual feature maps capturing both global and local information from 
the scene.

\paragraph{Step 4: Object detection with YOLOv9.}
YOLOv9 processes the extracted features and predicts bounding boxes for the objects 
of interest. In this application, the model is fine-tuned on the Michelin dataset 
to detect:
\begin{itemize}
    \item QR codes (used as calibration markers),
    \item the rubber strip,
    \item the cut edges of the rubber strip.
\end{itemize}
The output of this stage is a set of bounding boxes localizing the cut edges in 
pixel coordinates.

\paragraph{Step 5: Segmentation with EfficientSAM.}
The bounding boxes produced by YOLOv9 are passed as spatial prompts to EfficientSAM, 
which generates pixel-level segmentation masks for the detected cut edges. EfficientSAM 
processes each bounding box independently and outputs a binary mask indicating the 
exact pixel region corresponding to each cut edge.

\paragraph{Step 6: Edge extraction and post-processing.}
Once segmentation masks are obtained, classical image processing techniques are applied 
to extract geometric information:
\begin{itemize}
    \item contour detection,
    \item edge extraction,
    \item boundary analysis.
\end{itemize}
These operations allow identifying the precise boundary coordinates of each segmented 
cut edge.

\paragraph{Step 7: Geometric analysis and metric estimation.}
Using the segmented masks and the homography derived from the QR codes, the system 
computes:
\begin{itemize}
    \item the position of the cut edges in the coordinate system defined by the 
    QR codes, where point $(0, 0)$ is the centre of the upper-left QR code,
    \item the distance from each cut edge to the lower border of the table,
    \item the distance between the two cut edges.
\end{itemize}
All measurements are expressed in millimetres using the pixel-to-millimetre conversion 
factor estimated in the preprocessing step.

\subsubsection{Output}

The output of the YOLOv9 + EfficientSAM pipeline includes:
\begin{itemize}
    \item bounding boxes of the detected cut edges,
    \item pixel-level segmentation masks of the rubber strip cut edges,
    \item pixel coordinates of the detected edge boundaries,
    \item metric measurements of edge positions and distances in millimetres, 
    referenced to the QR code coordinate system.
\end{itemize}

\subsubsection{Evaluation metrics}

The performance of the pipeline is evaluated using the following metrics. For the 
object detection stage performed by YOLOv9: Average Precision (AP) and mean Average 
Precision (mAP), which evaluate the accuracy of predicted bounding boxes with respect 
to ground truth annotations. For the segmentation stage performed by EfficientSAM: 
Intersection over Union (IoU) and mask IoU, which measure the overlap between the 
predicted segmentation mask and the ground truth mask. These performance metrics are 
described in detail in Section~\ref{sec:metrics}.

\subsubsection{Applicability to the Michelin Challenge}

YOLOv9 + EfficientSAM is well suited for the Michelin Challenge 2 because it provides 
a complete and efficient pipeline for detection, segmentation and metric estimation. 
YOLOv9 delivers accurate and fast detection of cut edges even under variable 
illumination and perspective conditions, while EfficientSAM produces precise 
pixel-level segmentation masks at a fraction of the computational cost of the original 
SAM. The use of QR codes as calibration markers and the homography-based perspective 
correction ensure accurate metric estimation regardless of the camera position. 
Together, these components enable reliable localization and measurement of rubber 
strip cut edges in the automated inspection system.
% TODO: \subsection{Traditional Method (to be completed by group member)}

% --------------------
\section{Datasets}
% --------------------
\label{sec:datasets}
% At least two publicly available datasets per student, in paragraph format.
% Include name, date, link, and description for each dataset.

Several publicly available datasets are commonly used to evaluate object detection and 
segmentation models such as Grounding DINO and the Segment Anything Model. In 
addition, this project makes use of a proprietary dataset provided by Michelin 
Aranda de Duero as part of the challenge.

The \textbf{COCO (Common Objects in Context)} dataset was introduced in 2014 by 
Microsoft~\cite{LinCOCO2014} and is available at \url{https://cocodataset.org}. 
COCO contains more than 330,000 images, with over 200,000 labeled images and more 
than 1.5 million object instances annotated across 80 object categories, covering 
common objects in everyday scenes such as people, vehicles, animals, and household 
items. Each image provides detailed annotations including bounding boxes, instance 
segmentation masks, and keypoints, making the dataset suitable for evaluating object 
detection, instance segmentation, and keypoint detection methods. Due to its diversity 
and scale, COCO is the primary benchmark for evaluating modern detection models such 
as Grounding DINO and YOLOv8.Two samples from the dataset with their segmentation are shown in Figure~\ref{fig:coco_samples}\\

\begin{figure}[htb!]
    \centering
    \begin{subfigure}[b]{0.48\linewidth}
        \includegraphics[width=\linewidth, height=5cm, keepaspectratio]
        {figures/coco_times_original.jpg}
    \end{subfigure}
    \hfill
    \begin{subfigure}[b]{0.48\linewidth}
        \includegraphics[width=\linewidth, height=5cm, keepaspectratio]
        {figures/coco_times_segmented.jpg}
    \end{subfigure}
    \vspace{0.3cm}
    \begin{subfigure}[b]{0.48\linewidth}
        \includegraphics[width=\linewidth, height=5cm, keepaspectratio]
        {figures/coco_dog_original.jpg}
    \end{subfigure}
    \hfill
    \begin{subfigure}[b]{0.48\linewidth}
        \includegraphics[width=\linewidth, height=5cm, keepaspectratio]
        {figures/coco_dog_segmented.jpg}
    \end{subfigure}
    \caption{Sample images from the COCO dataset~\cite{LinCOCO2014}. 
    Left column: original images. Right column: instance segmentation 
    annotations showing object masks across multiple categories.}
    \label{fig:coco_samples}
\end{figure}

Another relevant dataset is \textbf{NEU Surface Defect Database} was introduced in 2013 by Northeastern 
University~\cite{NeuDataset} and is available at 
\url{https://ieee-dataport.org/documents/neu-det}. 
The dataset contains 1,800 grayscale images of hot-rolled steel strip surfaces, 
divided into six categories of surface defects: rolled-in scale, patches, crazing, 
pitted surface, inclusion, and scratches. Each category contains 300 samples, and 
the images have a resolution of $200 \times 200$ pixels. The dataset provides both 
image-level labels for classification tasks and bounding box annotations for object 
detection tasks, making it particularly suitable for evaluating defect detection 
methods in industrial inspection scenarios. Due to its focus on surface defect 
localization in manufacturing environments, NEU is directly relevant to the 
Michelin Challenge, where precise detection and localization of structural 
anomalies in rubber strips is required. Some images of the dataset are shown in Figure~\ref{fig:neu_samples}\\

\begin{figure}[htb!]
    \centering
    \begin{subfigure}[b]{0.48\linewidth}
        \includegraphics[width=\linewidth, height=5cm, keepaspectratio]
        {figures/neu_sample1.png}
    \end{subfigure}
    \hfill
    \begin{subfigure}[b]{0.48\linewidth}
        \includegraphics[width=\linewidth, height=5cm, keepaspectratio]
        {figures/neu_sample2.png}
    \end{subfigure}
    \caption{Sample images from the NEU Surface Defect Database~\cite{NeuDataset}, 
    showing bounding box annotations for surface defect detection on hot-rolled 
    steel strips.}
    \label{fig:neu_samples}
\end{figure}

The \textbf{LVIS (Large Vocabulary Instance Segmentation)} dataset was introduced 
in 2019 by Gupta et al. at Facebook AI Research and is available at its official 
website\footnote{\url{https://www.lvisdataset.org}}. LVIS contains over 164,000 
images sourced from the COCO image collection, annotated with high-quality instance 
segmentation masks across 1,203 object categories, making it one of the most 
diverse and fine-grained instance segmentation benchmarks available. The dataset 
is divided into three frequency-based splits: frequent categories (more than 100 
training samples), common categories (between 10 and 100 samples), and rare 
categories (fewer than 10 samples), reflecting the long-tail distribution of 
object categories in the real world. Annotations were collected using a 
federated annotation protocol that assigns only a subset of categories to each 
image, resulting in over 2 million high-quality segmentation masks. LVIS is used 
to evaluate EfficientSAM \cite{Xiong2023EfficientSAM} on zero-shot instance 
segmentation, where EfficientSAM-S achieves 42.3 AP on LVIS compared to 44.7 AP 
for the original SAM, while using $20\times$ fewer parameters.

%%%%%%%%%%%%%%%%%%%%%%%%%%%%%%%%%%%%%%%%%%%%%%%%%%%%%%%%%%%%%%%%%
%%%%%%%%%%%%%%%%%%%%%%%%%%%%%%%%%%%%%%%%%%%%%%%%%%%%%%%%%%%%%%%%%
%% FALTA AÑADIR LAS IMAGENES (nombres LVIS_...)
%%%%%%%%%%%%%%%%%%%%%%%%%%%%%%%%%%%%%%%%%%%%%%%%%%%%%%%%%%%%%%%%%
%%%%%%%%%%%%%%%%%%%%%%%%%%%%%%%%%%%%%%%%%%%%%%%%%%%%%%%%%%%%%%%%%

The \textbf{SA-1B} dataset was introduced in 2023 by Kirillov et al. at Meta FAIR 
and is available at its official 
website\footnote{\url{https://ai.meta.com/datasets/segment-anything/}}. SA-1B is 
the largest image segmentation dataset to date, comprising over 11 million 
high-resolution images with more than 1.1 billion automatically generated segmentation 
masks. The masks cover objects of all categories without semantic constraints, 
including whole objects and their parts and subparts, making it a general-purpose 
dataset for training and evaluating promptable segmentation models. It is used to 
pretrain and evaluate both SAM 2 \cite{Ravi2024SAM2} and EfficientSAM 
\cite{Xiong2023EfficientSAM}, with the latter achieving 44.4 AP on COCO zero-shot 
instance segmentation after fine-tuning on SA-1B.\\

Lastlly, we have \textbf{Michelin Aranda Dataset}, which is a propietary image collection provided by
Michelin Aranda. The dataset consists of images 
captured in laboratory conditions simulating real factory scenarios, using a flat table 
as a surface. Each image contains rubber strips placed on the table alongside three QR 
codes arranged in a right triangle with legs of 800 mm, used as calibration markers 
for pixel-to-metric conversion. The images were captured with a smartphone under 
variable camera positions and angles, reflecting the real acquisition conditions on 
the production floor. The dataset is specifically designed for the tasks of QR 
detection, rubber strip segmentation, cut edge localization, and metric distance 
estimation. Two sample images from this dataset are shown in 
Figure~\ref{fig:michelin_samples}.

\begin{figure}[htb!]
    \centering
    \begin{subfigure}[b]{0.48\linewidth}
        \includegraphics[width=\linewidth, height=6cm, keepaspectratio]{figures/michelin1.jpg}
    \end{subfigure}
    \hfill
    \begin{subfigure}[b]{0.48\linewidth}
        \includegraphics[width=\linewidth, height=6cm, keepaspectratio]{figures/michelin2.jpg}
    \end{subfigure}
    \caption{Sample images from the Michelin Aranda Dataset. QR codes are visible 
    in the corners of the surface and serve as spatial calibration markers.}
    \label{fig:michelin_samples}
\end{figure}

% --------------------
\section{Evaluation Metrics}
\label{sec:metrics}

The performance of the proposed system is evaluated using standard metrics commonly 
adopted in object detection and image segmentation tasks.

\paragraph{Precision.}
Precision measures the proportion of correctly predicted detections among all predicted 
detections. It reflects the model's ability to avoid false positives.

\[
\text{Precision} = \frac{TP}{TP + FP}
\]

\paragraph{Recall.}
Recall measures the ability of the model to detect all relevant objects present in the 
image. It reflects the model's capacity to minimize false negatives.

\[
\text{Recall} = \frac{TP}{TP + FN}
\]

\paragraph{Intersection over Union (IoU).}
Intersection over Union (IoU) evaluates the overlap between the predicted region and the 
ground truth annotation. It is defined as the ratio between the intersection area and the 
union area of the predicted and ground truth regions.

\[
\text{IoU} = \frac{\text{Area of Overlap}}{\text{Area of Union}}
\]

IoU is widely used in both object detection and segmentation tasks. A prediction is 
typically considered correct if the IoU exceeds a predefined threshold (e.g., 0.5), 
although more strict evaluations may use multiple thresholds.

\paragraph{Average Precision (AP).}
Average Precision summarizes the precision-recall curve for a given class by computing 
the area under the curve. It provides a single value representing the model performance 
across different recall levels.

\[
\text{AP} = \int_0^1 \text{Precision}(\text{Recall}) \, d(\text{Recall})
\]

\paragraph{Mean Average Precision (mAP).}
Mean Average Precision (mAP) is computed as the mean of the Average Precision across all 
object classes and is one of the most widely used metrics in object detection.

\[
\text{mAP} = \frac{1}{N} \sum_{i=1}^{N} AP_i
\]

where $N$ is the number of object classes and $AP_i$ is the Average Precision for class $i$.

\paragraph{Distance Error.}
In addition to detection and segmentation metrics, geometric accuracy is evaluated 
by measuring the error in the estimated distances between key elements, such as 
cutting edges and reference QR markers. This metric is defined as the absolute 
difference between the predicted distance and the ground truth distance in 
millimetres:

\[
\text{Distance Error} = | d_{\text{pred}} - d_{\text{gt}} |
\]

where $d_{\text{pred}}$ is the estimated distance and $d_{\text{gt}}$ is the 
ground truth measurement. This metric is particularly important in the context of 
the Michelin Challenge, where precise spatial measurements are required for 
industrial automation.\\

In the context of the Michelin Challenge, mAP is used to evaluate the detection 
of QR codes and relevant objects, IoU is particularly important for assessing the 
accuracy of segmentation masks and edge localization, and the Distance Error 
metric directly evaluates the precision of the geometric measurements required 
for industrial automation.

% --------------------
% --------------------
\section{Python Implementations}
% --------------------
\label{sec:python}
% Provide links to GitHub/GitLab repositories, release dates, advantages, 
% and limitations for each method. Include a comparison at the end.

This section presents the available Python implementations of the reviewed methods, 
including their repositories, release dates, main advantages, and limitations, with 
a focus on their practical applicability in real-world scenarios.

\textbf{YOLOv8} \cite{Jocher2023YOLOv8} is available at 
\url{https://github.com/ultralytics/ultralytics}. The official implementation was 
released on January 10, 2023, provided by Ultralytics as a unified Python package 
available via pip. It is built on PyTorch and offers both a Python API and a 
command-line interface, supporting multiple model sizes (n, s, m, l, x) and export 
to ONNX, TensorRT, and CoreML formats. Its main advantages are ease of use, 
real-time inference capability, and support for advanced data augmentation strategies 
such as mosaic and mixup. A possible limitation is that YOLOv8 operates on a closed 
vocabulary, making it less flexible than open-vocabulary methods when unseen 
categories need to be detected. Additionally, larger model variants require 
significantly more computational resources.

\textbf{RT-DETRv4} \cite{LiaoRTDETRv4} is available at 
\url{https://github.com/RT-DETRs/RT-DETRv4}. The repository was created on 
October 30, 2025, with full code, configurations, and checkpoints released on 
November 17, 2025. The implementation is provided in PyTorch and includes full 
training, evaluation, and deployment scripts for four model sizes (S, M, L, X) 
using HGNetv2 as backbone. Its main advantages include high detection accuracy 
through transformer-based global attention, improved efficiency over earlier 
DETR-based detectors, and an end-to-end pipeline that avoids Non-Maximum Suppression 
(NMS). Its main limitations are higher computational requirements than lightweight 
CNN detectors such as YOLOv8 and greater implementation complexity.

\textbf{AnoDDPM} \cite{WyattAnoDDPM2022} is available at 
\url{https://github.com/Julian-Wyatt/AnoDDPM}. The method was introduced at CVPR 
Workshops 2022. The implementation is written in PyTorch and provides scripts for 
training diffusion models on normal data and running anomaly detection at inference 
time, supporting both Gaussian and simplex noise schedules. Note that the repository 
was archived by the author in October 2025 and is no longer actively maintained. 
Its main advantages are the ability to learn from defect-free data only and stable 
training compared to GAN-based approaches. Its main limitations are high 
computational cost due to the iterative denoising process and the requirement for 
careful tuning of noise schedules.

\textbf{Segformer++} \cite{KienzleSegformer2024} is available at 
\url{https://github.com/KieDani/SegformerPlusPlus}. It was introduced in May 2024 
and is built on PyTorch with integration with the MMSegmentation and MMPose 
libraries. A key practical advantage is that the token merging strategy can be 
applied at inference time using the original Segformer pre-trained weights, without 
any retraining. Its main limitations are that performance gains are less pronounced 
at lower resolutions, and the method is specifically adapted to the Segformer 
architecture.

\textbf{PEM} \cite{CavagneroPEM2024} is available at 
\url{https://github.com/NiccoloCavagnero/PEM}. It was introduced on February 29, 
2024, and is provided in PyTorch with support for both semantic and panoptic 
segmentation under a single unified codebase. Its main advantages are significantly 
faster inference than comparable models such as Mask2Former and a unified framework 
that avoids task-specific architectures. Its main limitations are that performance 
falls slightly below the most accurate but slower models on some benchmarks, and 
computational requirements remain higher than lightweight CNN-only architectures.

\textbf{Relation-DETR} \cite{HouRelationDETR2024} is available at 
\url{https://github.com/xiuqhou/Relation-DETR}. It was introduced on July 16, 2024, 
and is provided in PyTorch with full training and evaluation scripts for COCO 2017 
and several domain-specific benchmarks. The position relation encoder is designed as 
a plug-in-and-play module compatible with any DETR-based architecture. Its main 
advantages are significantly faster convergence and state-of-the-art detection 
performance. Its main limitations are increased memory requirements at training time 
and limited validation on domain-specific datasets.

\textbf{ViLU-Net} \cite{HeidariViLUNet2025} is available at 
\url{https://github.com/moeinheidari7829/Retroperitoneal_Tumour_Segmentation_SPIE2025}. 
It was introduced on February 1, 2025, and is built on top of the nnU-Net framework 
in PyTorch. Its main advantages are state-of-the-art performance on medical 
segmentation benchmarks and efficient long-range dependency modeling via xLSTM blocks. 
Its main limitations are the small in-house dataset used for evaluation, the high GPU 
memory requirements (NVIDIA V100, 32 GB), and limited validation beyond CT imaging.

\textbf{Grounded-SAM} \cite{RenGroundedSAM2024} is available at 
\url{https://github.com/IDEA-Research/Grounded-Segment-Anything}. It was introduced 
on January 25, 2024, and is actively maintained. The implementation follows a 
training-free pipeline combining Grounding DINO with SAM, and includes ready-to-use 
scripts for multiple pipelines including RAM-Grounded-SAM and integration with 
HQ-SAM and faster SAM variants such as MobileSAM and EfficientSAM. Its main 
advantages are open-vocabulary detection and segmentation without retraining and 
strong zero-shot performance on the SegInW benchmark. Its main limitations are 
inference speed bottlenecks for real-time applications and dependence on the quality 
of Grounding DINO detections.

\textbf{YOLO-World} \cite{ChengYOLOWorld2024} is available at 
\url{https://github.com/AILab-CVC/YOLO-World}. It was introduced on January 30, 
2024, and is built on top of the MMYOLO and MMDetection toolboxes in PyTorch. It 
provides three model variants (S, M, L) and supports pre-training on large-scale 
datasets as well as fine-tuning for downstream tasks. Its main advantages are 
real-time open-vocabulary detection at 52 FPS and a lightweight architecture with 
as few as 13M parameters. Its main limitations are lower zero-shot performance than 
heavier open-vocabulary detectors and high pre-training computational requirements.

\textbf{Grounding DINO 1.5} \cite{RenGroundingDINO15} is available at 
\url{https://github.com/IDEA-Research/Grounding-DINO-1.5-API}. It was introduced on 
May 16, 2024, and is accessible via API rather than as a fully open-source codebase. 
It includes two model variants: the Pro model using a ViT-L backbone for maximum 
accuracy, and the Edge model using EfficientViT-L1 for real-time inference. Its main 
advantages are state-of-the-art zero-shot detection performance and support for 
diverse input modalities. Its main limitations are the closed API restricting 
reproducibility and customization, and the high computational requirements of the 
Pro model.

The reviewed methods can be grouped into three main categories: object detection, 
segmentation, and anomaly detection. Object detection methods --- YOLOv8, RT-DETRv4, 
Relation-DETR, YOLO-World, Grounding DINO 1.5, and Grounded-SAM --- differ in their 
design philosophy. YOLO-based approaches provide efficient real-time performance, 
while transformer-based detectors such as RT-DETRv4 and Relation-DETR achieve higher 
accuracy through end-to-end architectures without the need for Non-Maximum Suppression. 
Open-vocabulary models such as Grounding DINO 1.5 and Grounded-SAM further extend 
these capabilities by enabling flexible detection beyond predefined categories.

Segmentation methods --- Segformer++, PEM, and ViLU-Net --- follow encoder-decoder 
architectures for dense prediction tasks, differing mainly in their efficiency 
strategies and ability to balance accuracy and computational cost. In contrast, 
AnoDDPM represents a fundamentally different paradigm, as an unsupervised generative 
approach that learns normal data distributions and detects anomalies without requiring 
labeled defect samples.

From the perspective of the Michelin Challenge, methods such as Grounded-SAM and 
YOLO-World are particularly suitable due to their open-vocabulary capabilities and 
generalization to unseen scenarios. Additionally, YOLO-based approaches offer a strong 
advantage in terms of real-time performance, which is critical in industrial inspection 
settings. However, diffusion-based methods such as AnoDDPM, while effective in 
low-supervision scenarios, may be limited by their high computational cost in 
real-time deployment environments.

%%%%%%%%%%%%%%%%%%%%%%%%%%%%%%%%%%%%%%%%%%%%%%%%
%%%%%%%%%%%%%%%%%%%%%%%%%%%%%%%%%%%%%%%%%%%%%%%%
%% IMPLEMENTACIONES JAIME
%%%%%%%%%%%%%%%%%%%%%%%%%%%%%%%%%%%%%%%%%%%%%%%%
%%%%%%%%%%%%%%%%%%%%%%%%%%%%%%%%%%%%%%%%%%%%%%%%

\textbf{SAM 2} \cite{Ravi2024SAM2} is available in its official 
repository\footnote{\url{https://github.com/facebookresearch/sam2}}. The official 
implementation was released on August 1, 2024, by Meta FAIR, provided as an 
open-source PyTorch package under an Apache 2.0 license. It supports both image 
and video segmentation through a unified promptable interface accepting points, 
bounding boxes, and masks as inputs. Multiple model sizes are available (Hiera-T, 
S, B+, L) with the B+ variant running at 43.8 FPS on a single A100 GPU. Its main 
advantages are strong zero-shot segmentation performance without task-specific 
training, a streaming memory architecture enabling video processing, and active 
maintenance from Meta. A possible limitation is that larger model variants require 
significant GPU memory, and the model produces class-agnostic masks that require 
additional classification logic for specific downstream tasks.\\

\textbf{EfficientSAM} \cite{Xiong2023EfficientSAM} is available in its official 
repository\footnote{\url{https://github.com/yformer/EfficientSAM}}. The 
implementation was released in December 2023 by Meta AI Research in PyTorch. It 
provides two model variants, EfficientSAM-Ti (10M parameters) and EfficientSAM-S 
(25M parameters), both pretrained via SAMI on SA-1B and fine-tuned for the segment 
anything task. Its main advantages are a $20\times$ reduction in parameter count 
and inference time compared to SAM ViT-H while maintaining competitive segmentation 
quality, with EfficientSAM-S achieving 44.4 AP on COCO zero-shot instance 
segmentation versus 46.5 AP for SAM. Its main limitations are reduced documentation 
compared to the SAM 2 repository, the absence of video segmentation support, and 
lower performance on fine-grained segmentation tasks compared to SAM 2.\\

\textbf{YOLOv9} \cite{Wang2024YOLOv9} is available in its official 
repository\footnote{\url{https://github.com/WongKinYiu/yolov9}}. The implementation 
was released on February 29, 2024, by the original authors in PyTorch with full 
training, evaluation, and inference scripts for MS COCO. It provides multiple model 
sizes (S, M, C, E) and supports both object detection and instance segmentation 
tasks. Its main advantages are state-of-the-art accuracy among train-from-scratch 
detectors, achieving 55.6\% AP$_{50:95}$ with YOLOv9-E while reducing parameters 
by 49\% and computation by 43\% compared to YOLOv8-X, and well-documented training 
pipelines that facilitate fine-tuning on custom datasets. Its main limitations are 
that the PGI auxiliary branch adds training complexity, and the repository has seen 
reduced maintenance activity compared to the Ultralytics ecosystem.\\

\textbf{YOLOv8-HDSA} \cite{Shi2025CarbonBlock} is based on the Ultralytics YOLOv8 
framework, available in its official 
repository\footnote{\url{https://github.com/ultralytics/ultralytics}}. The base 
framework was released in January 2023 and is actively maintained by Ultralytics 
under an AGPL-3.0 license, with the HDSA modifications introduced in the paper 
published in March 2025. The Ultralytics package provides a highly accessible Python 
API and command-line interface supporting detection, segmentation, classification, 
and pose estimation across multiple model sizes. Its main advantages are an 
exceptionally mature ecosystem with extensive documentation, active community 
support, and easy fine-tuning on custom datasets via a single configuration file. 
Its main limitation for this specific paper is that the HDSA-specific modifications 
(SRFF module, convolutional self-attention head, Focaler-IoU loss) are not released 
separately and must be reimplemented from the paper description.\\

\textbf{AnomalyDINO} \cite{Damm2024AnomalyDINO} is available in its official 
repository\footnote{\url{https://github.com/SimonDamm/AnomalyDINO}}. The 
implementation was released in May 2024 alongside the paper submission and accepted 
at WACV 2025. It is provided in PyTorch and leverages DINOv2 pretrained weights 
from the \texttt{torch.hub} registry, requiring no additional training data beyond 
a few reference normal images. Its main advantages are a completely training-free 
inference pipeline, support for both image-level anomaly scoring and pixel-level 
anomaly segmentation, and state-of-the-art one-shot performance achieving 96.6\% 
AUROC on MVTec-AD. Its main limitations are that performance degrades when the 
number of reference images is very small (below 3--5), and the method is sensitive 
to domain shift between reference and test images.\\

\textbf{CLIP-DINOv2} \cite{Jiang2025CLIPDINO} does not have a dedicated public 
repository. The paper, published in Electronics in December 2025, provides full 
implementation details including architecture specifications, hyperparameters, and 
training procedures. The method relies on publicly available pretrained weights for 
CLIP ViT-L/14@336px, available via 
OpenAI\footnote{\url{https://github.com/openai/CLIP}}, and DINOv2 ViT-L, available 
via \texttt{torch.hub}\footnote{\url{https://github.com/facebookresearch/dinov2}}. 
Its main advantages are state-of-the-art zero-shot anomaly detection performance 
across seven industrial benchmarks without requiring target-domain samples. Its main 
limitation is the absence of a public codebase, which requires full reimplementation 
from the paper to reproduce results.\\

\textbf{SAM2-UNet} \cite{Xiong2026SAM2UNet} is available in its official 
repository\footnote{\url{https://github.com/WZH0120/SAM2-UNet}}. The preprint was 
released in 2024 and the peer-reviewed version was published in Visual Intelligence 
in January 2026. The implementation is provided in PyTorch and requires SAM2 
pretrained Hiera weights available from the SAM 2 
repository\footnote{\url{https://github.com/facebookresearch/sam2}}. It supports 
five benchmarks out of the box with ready-to-use training and evaluation scripts. 
Its main advantages are parameter-efficient fine-tuning via lightweight adapters 
that freeze the Hiera encoder, enabling training on a single 24 GB GPU, and strong 
generalization across diverse segmentation tasks without task-specific architectural 
changes. Its main limitation is that the current implementation operates at 
352$\times$352 resolution, which may reduce performance on high-resolution images 
with fine edge details.\\

\textbf{Semi-supervised defect segmentation} \cite{Shi2024SemiSupervised} does not 
have a dedicated public repository. The paper was published in Scientific Reports in 
September 2024, with code available upon request from the corresponding 
author\footnote{\url{https://doi.org/10.1038/s41598-024-72579-6}}. The implementation 
is based on PyTorch with a modified DeepLabv3Plus backbone using a simplified 
MobileNetv2 encoder. Its main advantages are a lightweight architecture achieving 
64.59 FPS on an RTX 3090 while outperforming SOTA semi-supervised methods, and a 
dynamic threshold strategy that improves pseudo-label quality over training. Its 
main limitation is the lack of a public repository, which hinders reproducibility 
and requires direct contact with the authors.\\

\textbf{Camera Calibration Survey} \cite{Liao2025CameraCalibration} maintains a 
curated repository\footnote{\url{https://github.com/KangLiao929/Awesome-Deep-Camera-Calibration}}. 
The survey was originally submitted in March 2023 and last revised in February 2025 
(v3). The repository is not a single codebase but a structured collection of links 
to implementations of calibration methods organized by category (pinhole, distortion, 
cross-view, cross-sensor), alongside a unified holistic calibration dataset for 
benchmarking. Its main advantage is providing a comprehensive reference point for 
selecting and comparing calibration methods relevant to the QR code-based spatial 
calibration required in the Michelin Challenge. Its main limitation is that it does 
not provide ready-to-use calibration code and requires selecting and integrating 
individual method repositories.\\

\textbf{MVTec AD 2} \cite{HecklerKram2025MVTecAD2} provides the dataset and 
evaluation scripts at its official 
website\footnote{\url{https://www.mvtec.com/company/research/datasets/mvtec-ad-2}}. 
The dataset was released in March 2025 by MVTec Software GmbH. It does not provide 
a segmentation or detection model but instead offers a standardized evaluation 
framework compatible with existing anomaly detection pipelines. Its main advantage 
is providing a challenging and realistic benchmark with scenarios not covered by 
existing datasets, including transparent objects, dark-field illumination, and 
extremely small defects. Its main limitation is that no baseline model is included, 
requiring users to integrate their own anomaly detection methods for evaluation.\\

The reviewed methods can be grouped into three main categories: foundation 
segmentation models, industrial defect detection and segmentation, and supporting 
tools. Foundation models --- SAM 2, EfficientSAM, and SAM2-UNet --- share a common 
architectural lineage based on the Hiera encoder and promptable segmentation paradigm, 
differing mainly in computational efficiency and downstream adaptability. SAM 2 
offers the strongest zero-shot performance, EfficientSAM prioritizes inference speed, 
and SAM2-UNet enables supervised fine-tuning through a U-shaped decoder without 
modifying the encoder.

Industrial defect methods --- YOLOv9, YOLOv8-HDSA, AnomalyDINO, CLIP-DINOv2, and 
the semi-supervised segmentation approach --- differ fundamentally in their supervision 
requirements. YOLOv9 and YOLOv8-HDSA require labeled bounding box and mask 
annotations for fine-tuning, while AnomalyDINO and CLIP-DINOv2 operate in zero-shot 
and few-shot settings without target-domain training data. The semi-supervised method 
occupies an intermediate position, exploiting both labeled and unlabeled samples to 
maximize data utilization in annotation-scarce scenarios.

From the perspective of the Michelin Challenge, where the available dataset is 
small and annotations are limited, the most practical approaches are AnomalyDINO 
for anomaly detection due to its training-free nature, and the YOLOv9 + EfficientSAM 
pipeline for detection and segmentation, as both are publicly available, 
well-documented, and compatible with PyTorch-based fine-tuning on small custom datasets.
% --------------------
\section{CREDITS}
% --------------------
\label{sec:credits}
Review the following link for your information: \footnote{https://www.elsevier.com/researcher/author/policies-and-guidelines/credit-author-statement} “CRediT (Contributor Roles Taxonomy) was introduced to recognize individual author contributions, reducing authorship disputes and facilitating collaboration”. More informally, it is a paragraph usually introduced at the end of research papers that recognizes the individual work of each author. You will see this in most of the papers you will be reviewing. 

The group will include a paragraph like the CRediT statement indicating who collaborated on each section and task of the research proposal. Not all the original CRediT Terms would be applicable, so the group could use the following:

\begin{table}[ht]
\centering
\begin{tabular}{|l|p{10cm}|}
\hline
\textbf{Term} & \textbf{Definition} \\ \hline
Investigation & Conducting a research and investigation process \\ \hline
Resources & Collected or gathered datasets mentioned in the study \\ \hline
Writing - Original Draft & Preparation of the research proposal, explicitly writing the initial draft. \\ \hline
Writing - Review \& Editing & Critical review, commentary, or revision \\ \hline
Visualization & Preparation of visual data, including figures and summary tables \\ \hline
Supervision & Oversight and leadership responsibility for the research activity planning and execution, including external mentorship to the core team \\ \hline
Project POC & Responsibility for being the point of contact with the professor for doubts and task upload. \\ \hline
\end{tabular}
\caption{Definitions of Research Terms}
\label{tab:definitions}
\end{table}


For a group of three people, this could be a template.
\begin{itemize}
\item N\_Surname\_01 contributed to the Investigation of XXXX methods and participated in Writing - Original Draft, preparing the initial manuscript. Additionally, they served as the Project POC, acting as the main point of contact for the professor to address any doubts or concerns.

\item N\_Surname\_02 also participated in the Investigation of YYYY methods and Writing - Original Draft, contributing to the research process and drafting the manuscript. They were responsible for resources, collecting and organizing datasets, and studying materials essential for the project.

\item N\_Surname\_03 equally contributed to the Investigation of ZZZZ methods and Writing - Original Draft, engaging in the research and drafting activities. They took charge of Visualization, preparing figures and summary tables to present the study findings effectively, and provided Supervision, ensuring oversight and leadership throughout the research activities.
 
\end{itemize}


% %%%%%%%%%%%%%%%%%%%%%%%%%%%%%%%%%%%%%%%%%%%%%%%%%%%%%%%%%%
% %%%%%%%%%%%%%%%%%%%%%%%%%%%%%%%%%%%%%%%%%%%%%%%%%%%%%%%%%%
% REFERENCES SECTION
% %%%%%%%%%%%%%%%%%%%%%%%%%%%%%%%%%%%%%%%%%%%%%%%%%%%%%%%%%%
% %%%%%%%%%%%%%%%%%%%%%%%%%%%%%%%%%%%%%%%%%%%%%%%%%%%%%%%%%%
\medskip

\bibliography{references.bib} 

\newpage


\end{document}